\chapter{Polinomios}

\section{Definiciones y operaciones}
\begin{definicion}
Sea $K$ un cuerpo, por ejemplo $\mathbb Q$, $\mathbb R$, $\mathbb C$
o $\mathbb F_p=\mathbb Z/p\mathbb Z$, con $p$ primo.
Un \textbf{polinomio} de $K[X]$ es una expresión finita
\[
f=a_0+a_1X+\cdots+a_nX^n,\qquad a_i\in K.
\]
Si $a_n\neq0$, su \textbf{grado} es $\deg f=n$ y su
\textbf{coeficiente principal} es $a_n$. Es \textbf{mónico} si $a_n=1$.
El polinomio cero no tiene grado en la convención de este apunte.
Dos polinomios son iguales si coinciden todos sus coeficientes.
\end{definicion}
\begin{proposicion}[{title={Operaciones y grados}}]
La suma se hace coeficiente a coeficiente. En el producto,
\[
[X^k](fg)=\sum_{i+j=k}a_i b_j,
\]
donde $[X^k]f$ significa el coeficiente de $X^k$ en $f$.
Para $f,g\neq0$:
\[
\deg(fg)=\deg f+\deg g.
\]
Si $f+g\neq0$, $\deg(f+g)\leq\max\{\deg f,\deg g\}$;
hay igualdad si los grados son distintos. Si $g$ no es constante y
$f\neq0$ tiene grado $n$, entonces $\deg(f\circ g)=n\deg g$.
\end{proposicion}
\begin{ejemplo}
En $\mathbb Q[X]$, $\deg((X^2+1)^3)=6$.
El coeficiente de $X^4$ en $(1+X)^6$ es $\binom64=15$.
En $\mathbb F_2[X]$, $X+X=0$: reducir los coeficientes puede cambiar el grado.
\end{ejemplo}
\begin{nota}{Polinomio y función polinómica.}
Evaluar en $\alpha\in K$ significa sustituir $X$ por $\alpha$.
Sobre un cuerpo finito, dos polinomios distintos pueden dar la misma función:
$X^p-X$ no es el polinomio cero, aunque se anula en todo $\mathbb F_p$.
\end{nota}

\section{División, divisibilidad y restos}
\begin{proposicion}[{title={Algoritmo de división}}]
Dados $f,g\in K[X]$ con $g\neq0$, existen únicos $q,r\in K[X]$ tales que
\[
f=qg+r,\qquad r=0\ \text{o}\ \deg r<\deg g.
\]
Se dice que $g$ \textbf{divide} a $f$, y se escribe $g\mid f$, si $r=0$.
\end{proposicion}
\begin{nota}{Cómo dividir.}
Dividir el término de mayor grado del dividendo por el de $g$,
añadirlo al cociente y restar su producto por $g$.
Repetir hasta que el resto sea cero o tenga grado menor que $g$.
Los divisores constantes no nulos son unidades: dividen a todo polinomio.
\end{nota}
\begin{proposicion}[{title={Teorema del resto y de la raíz}}]
El resto de dividir $f$ por $X-\alpha$ es $f(\alpha)$. En particular,
\[
f(\alpha)=0\iff (X-\alpha)\mid f.
\]
Un polinomio no nulo de grado $n$ tiene a lo sumo $n$ raíces distintas en $K$.
\end{proposicion}
\begin{demostracion}
En $f=(X-\alpha)q+r$, el resto es constante. Evaluando en $\alpha$,
$r=f(\alpha)$. Cada raíz distinta aporta un factor lineal distinto;
la suma de sus grados no puede superar $\deg f$.
\end{demostracion}
\begin{proposicion}[{title={Identidades útiles}}]
Para $n\geq1$,
\[
X^n-a^n=(X-a)\sum_{j=0}^{n-1}X^{n-1-j}a^j.
\]
Sustituyendo $a$ por $-a$, se obtiene la factorización por $X+a$
de $X^n+a^n$ si $n$ es impar, y de $X^n-a^n$ si $n$ es par.
\end{proposicion}

\subsection{Congruencias polinómicas}
\begin{definicion}
Si $h\neq0$, se define $f\equiv g\pmod h$ cuando $h\mid f-g$.
Es una relación de equivalencia compatible con sumas, productos y
potencias de exponente entero no negativo.
El resto de $f$ módulo $h$ es el único representante $r$ con
$r=0$ o $\deg r<\deg h$.
\end{definicion}
\begin{ejemplo}
Módulo $X^2+1$ se tiene $X^2\equiv-1$ y $X^4\equiv1$.
Por tanto $X^{2026}+X^3\equiv-1-X$; ése es el resto.
\end{ejemplo}
\begin{nota}{Restos por evaluación.}
Si el divisor tiene grado $d$, escribir $r$ con grado menor que $d$.
En cada raíz $\alpha$ del divisor se cumple $r(\alpha)=f(\alpha)$.
Si hay $d$ raíces distintas, esas condiciones determinan el resto.
Si hay raíces repetidas, también se necesitan condiciones sobre derivadas.
Por ejemplo, sobre $\mathbb Q$, $\mathbb R$ o $\mathbb C$, el resto módulo
$(X-\alpha)^m$ es
\[
r=\sum_{j=0}^{m-1}\frac{f^{(j)}(\alpha)}{j!}(X-\alpha)^j.
\]
\end{nota}
\begin{ejemplo}
Si $f(1)=2$ y $f(-1)=4$, el resto módulo $X^2-1$ es $r=aX+b$.
De $a+b=2$, $-a+b=4$ resulta $r=-X+3$.
\end{ejemplo}

\section{Máximo común divisor y Bézout}
\begin{definicion}
Para $f,g\in K[X]$ no ambos nulos, su \textbf{mcd}, denotado $(f:g)$,
es el divisor común mónico de mayor grado. Son \textbf{coprimos} si $(f:g)=1$.
\end{definicion}
\begin{proposicion}
Si $f=qg+r$, entonces $(f:g)=(g:r)$.
El algoritmo de Euclides calcula el mcd mediante divisiones sucesivas:
es el último resto no nulo dividido por su coeficiente principal.
Además existen $A,B\in K[X]$ tales que
\[
(f:g)=Af+Bg.
\]
En particular, $(f:g)=1$ si y sólo si existe una combinación $Af+Bg=1$.
\end{proposicion}
\begin{nota}{Raíces comunes.}
Para $\alpha\in K$,
\[
f(\alpha)=g(\alpha)=0\iff(f:g)(\alpha)=0.
\]
Una dirección usa que el mcd divide a ambos; la otra se obtiene evaluando
la identidad de Bézout.
Si $f,g\in\mathbb C[X]$, ser coprimos equivale a no tener raíces comunes.
Sobre otros cuerpos, no basta con buscar raíces sólo en ese cuerpo.
\end{nota}
\begin{ejemplo}
Para $f=X^3-X$, $g=X^2-1$, se tiene $f=Xg$, luego $(f:g)=X^2-1$.
Las raíces comunes son $1$ y $-1$.
\end{ejemplo}

\section{Derivadas y multiplicidad}
\begin{definicion}
La \textbf{derivada formal} de $f=\sum_{j=0}^n a_jX^j$ es
$f'=\sum_{j=1}^n ja_jX^{j-1}$. Satisface las reglas usuales de suma
y producto. Si $f\neq0$, $\alpha$ tiene \textbf{multiplicidad} $m\geq1$ si
\[
f=(X-\alpha)^m q,\qquad q(\alpha)\neq0.
\]
Es simple si $m=1$ y múltiple si $m\geq2$.
\end{definicion}
\begin{proposicion}[{title={Criterios de multiplicidad}}]
Para $f\in K[X]$ no nulo, $\alpha$ es raíz múltiple si y sólo si
$f(\alpha)=f'(\alpha)=0$.
Sobre $\mathbb Q$, $\mathbb R$ o $\mathbb C$, tiene multiplicidad exacta $m$ si
\[
f(\alpha)=f'(\alpha)=\cdots=f^{(m-1)}(\alpha)=0,
\qquad f^{(m)}(\alpha)\neq0.
\]
En estos cuerpos, una raíz de multiplicidad $m$ en $f$ tiene multiplicidad
$m-1$ en $f'$ (para $m=1$, deja de ser raíz).
\end{proposicion}
\begin{demostracion}
Derivando $f=(X-\alpha)^m q$ se obtiene
\[
f'=(X-\alpha)^{m-1}\bigl(mq+(X-\alpha)q'\bigr).
\]
En característica cero el factor entre paréntesis vale $mq(\alpha)\neq0$
en $\alpha$. Derivando sucesivamente se obtiene el criterio completo.
En cualquier característica, si $f=(X-\alpha)q$, se tiene
$f'(\alpha)=q(\alpha)$, lo que prueba el criterio de raíz múltiple.
\end{demostracion}
\begin{nota}{Usar el mcd con la derivada.}
En $\mathbb C[X]$, todas las raíces de $f$ son simples si y sólo si
$(f:f')=1$. Si $f=c\prod_j(X-\alpha_j)^{m_j}$ con raíces distintas,
\[
(f:f')=\prod_j(X-\alpha_j)^{m_j-1}.
\]
Para buscar parámetros que produzcan raíces múltiples, resolver primero
$f'(\alpha)=0$ y luego imponer $f(\alpha)=0$.
En característica $p$ hay precauciones: $(X^p)'=0$; no aplicar sin revisar
las reglas de multiplicidad que requieren característica cero.
\end{nota}
\begin{ejemplo}
Sea $f=X^3-3X+a$. Como $f'=3(X^2-1)$, una raíz múltiple debe ser
$1$ o $-1$. Se cumple $f(1)=a-2$ y $f(-1)=a+2$.
Por tanto hay raíces múltiples exactamente para $a=2$ o $a=-2$.
Son dobles, pues $f''(\pm1)\neq0$.
\end{ejemplo}

\section{Raíces y factorización según el cuerpo}
\begin{proposicion}[{title={Raíces racionales: lema de Gauss}}]
Sea $f=a_nX^n+\cdots+a_0\in\mathbb Z[X]$, con $a_na_0\neq0$.
Si $u/v$, con $(u:v)=1$ y $v>0$, es raíz de $f$, entonces
\[
u\mid a_0,\qquad v\mid a_n.
\]
Para coeficientes racionales, primero multiplicar por un denominador común.
Si $a_0=0$, extraer la potencia de $X$ antes de buscar otras raíces.
\end{proposicion}
\begin{ejemplo}
En $f=2X^3-X^2-2X+1$, los candidatos racionales son
$\pm1,\pm\tfrac12$. Agrupando,
\[
f=(2X-1)(X^2-1)=(2X-1)(X-1)(X+1).
\]
Los candidatos del lema deben probarse: no todos son necesariamente raíces.
\end{ejemplo}
\begin{proposicion}[{title={Factorización en $\mathbb C$, $\mathbb R$ y $\mathbb Q$}}]
\begin{itemize}
\item En $\mathbb C[X]$, todo polinomio no constante tiene una raíz
(teorema fundamental del álgebra). Por tanto se factoriza completamente
en factores lineales; hay $\deg f$ raíces contando multiplicidades.
\item Si $f\in\mathbb R[X]$, las raíces no reales aparecen en pares conjugados
con la misma multiplicidad. El par $a\pm bi$, $b\neq0$, aporta el factor
$X^2-2aX+a^2+b^2$. Los irreducibles reales son los lineales y los
cuadráticos con discriminante negativo.
\item En $\mathbb Q[X]$, hay irreducibles de grados mayores. Si $\deg f=2$ o $3$,
es irreducible si y sólo si no tiene raíces racionales.
Para grado $4$ o mayor, no tener raíces racionales no es suficiente.
\end{itemize}
\end{proposicion}
\begin{ejemplo}
El polinomio $X^4-1$ se factoriza en $\mathbb Q[X]$ y $\mathbb R[X]$ como
$(X-1)(X+1)(X^2+1)$; en $\mathbb C[X]$, el último factor es $(X-i)(X+i)$.
Si $\alpha$ es una raíz de un polinomio racional, también es raíz
$\overline\alpha$. Además, si $\sqrt2$ es raíz, $-\sqrt2$ también lo es:
al evaluar, separar la parte racional y el coeficiente de $\sqrt2$.
\end{ejemplo}
\begin{nota}{Cambios de variable y raíces de la unidad.}
Si sólo aparecen potencias múltiplos de $d$, usar $Y=X^d$.
Por ejemplo $X^6+X^3-2=(X^3-1)(X^3+2)$:
se resuelven $X^3=1$ y $X^3=-2$ con la fórmula de raíces complejas.
También
\[
1+X+\cdots+X^n=\frac{X^{n+1}-1}{X-1}.
\]
Sobre $\mathbb C$, sus raíces son las de $G_{n+1}$ salvo $1$, todas simples.
Si $f\in\mathbb C[X]$ tiene una raíz real y $f=u+iv$ con $u,v\in\mathbb R[X]$,
esa raíz debe anular simultáneamente a $u$ y $v$: buscar su mcd.
\end{nota}

\subsection{Identidades notables: factorización de polinomios}

\begin{proposicion}
    \begin{itemize}
        \item \textbf{Diferencia de cuadrados:}
        \[
            a^2-b^2=(a-b)(a+b)
        \]

        \item \textbf{Suma de cuadrados:}
        No se factoriza en general en $\mathbb{R}$,
        pero sí en $\mathbb{C}$:
        \[
            a^2+b^2=(a-bi)(a+bi)
        \]

        \item \textbf{Diferencia de cubos:}
        \[
            a^3-b^3=(a-b)(a^2+ab+b^2)
        \]

        \item \textbf{Suma de cubos:}
        \[
            a^3+b^3=(a+b)(a^2-ab+b^2)
        \]
    \end{itemize}
\end{proposicion}


\section{Irreducibilidad y cuerpos finitos}
\begin{definicion}
Un polinomio $f\in K[X]$ de grado positivo es \textbf{irreducible}
si no se puede escribir como producto de dos polinomios de grados positivos.
Todo polinomio no nulo se factoriza como una constante no nula por
un producto de irreducibles mónicos, de manera única salvo el orden.
\end{definicion}
\begin{nota}{Criterios prácticos.}
En cualquier cuerpo, un polinomio de grado $2$ o $3$ es irreducible
si y sólo si no tiene raíces en ese cuerpo.
Para grado $4$, además de descartar raíces, hay que descartar divisores
cuadráticos irreducibles. En $\mathbb F_2[X]$, el único cuadrático
mónico irreducible es $X^2+X+1$.
No confundir un polinomio sin raíces en $K$ con uno irreducible de grado arbitrario.
\end{nota}
\begin{ejemplo}
En $\mathbb F_2[X]$, $X^4+X+1$ no tiene raíces, y su resto al dividir
por $X^2+X+1$ es $1$: es irreducible.
En cambio $X^4+X^2+1=(X^2+X+1)^2$ es reducible y tampoco tiene raíces en $\mathbb F_2$.
\end{ejemplo}
\begin{proposicion}[{title={Conteo de cuadráticos mónicos}}]
En $\mathbb F_p[X]$ hay $p^2$ polinomios mónicos de grado $2$.
De ellos, $p(p+1)/2$ son reducibles y $p(p-1)/2$ son irreducibles.
\end{proposicion}
\begin{demostracion}
Hay $p$ opciones para cada coeficiente de $X^2+aX+b$.
Los reducibles corresponden a pares de raíces no ordenados:
$\binom p2$ pares distintos y $p$ pares repetidos.
Los restantes son irreducibles.
\end{demostracion}

\section{Relaciones entre raíces y coeficientes}
\begin{proposicion}[{title={Fórmulas de Viète}}]
Si $f=a_n\prod_{j=1}^n(X-\alpha_j)$, contando multiplicidades, entonces
\[
\sum_j\alpha_j=-\frac{a_{n-1}}{a_n},\qquad
\sum_{i<j}\alpha_i\alpha_j=\frac{a_{n-2}}{a_n}\quad(n\geq2),\qquad
\prod_j\alpha_j=(-1)^n\frac{a_0}{a_n}.
\]
Se obtienen expandiendo el producto e igualando coeficientes.
\end{proposicion}
\begin{ejemplo}
Si $a,b,c$ son las raíces de $2X^3-3X^2+4X+1$, entonces
\[
a+b+c=\frac32,\qquad ab+ac+bc=2,\qquad abc=-\frac12.
\]
No hace falta hallar las raíces por separado.
\end{ejemplo}

\section{Cómo elegir el método}
\begin{nota}{Guía breve.}
\begin{enumerate}
\item \textbf{Grado o coeficiente:} identificar términos principales y
sumar sólo los productos que aportan a la potencia buscada.
\item \textbf{Resto:} usar congruencias polinómicas o escribir un resto
de grado menor que el divisor e imponer evaluaciones y, si hace falta, derivadas.
\item \textbf{Raíces comunes:} calcular el mcd por Euclides.
\item \textbf{Multiplicidad o parámetros:} combinar $f(\alpha)=0$ con
$f'(\alpha)=0$; seguir derivando para conocer la multiplicidad exacta.
\item \textbf{Factorización:} fijar el cuerpo, extraer factores evidentes,
buscar raíces racionales o usar un cambio de variable; conservar el coeficiente principal.
\item \textbf{Grado mínimo:} reunir todas las raíces forzadas, sus conjugadas
y sus multiplicidades; después ajustar la constante con las evaluaciones dadas.
\item \textbf{Sucesiones de polinomios:} usar inducción, comparando grados,
coeficientes o derivadas según la recurrencia.
\item \textbf{Sumas y productos de raíces:} usar Viète antes de intentar resolver el polinomio.
\end{enumerate}
\end{nota}
